\scanpage{60}
% archive foliation: 59
\deleted{
$C=(C_0,C_1)$ groupoïde connexe

$G=\operatorname{Aut}(C)$ gr. des automorphismes

$G^0$ ss gr. des $g\in G$ induisant l'identité sur $C_0$
\diagram{\begin{tikzcd}[ampersand replacement=\&]
1 \arrow[r] \& G^0 \arrow[r] \& G \arrow[r,"{\uuline{\text{Surj}}}"'] \& \unsure{\mathfrak S|\mathcal S}_{C_0} \arrow[r] \& 1
\end{tikzcd}}
[Pour qu'une autoéquivalence \inserted{$u$} soit un iso, il f et s. qu'elle soit bijective sur les objets de $C$, et induise un iso. \deleted{sur les} $\operatorname{Aut}(x)\longrightarrow\operatorname{Aut}(ux)$ pour $\forall x\in C_0$ ; d'autre part, pour toute application $C_0\overset{u_0}{\unsure{-|\longrightarrow}}C'_0$
}
\diagram{\ensuremath{\begin{array}{ccccccccc}
1&-&\operatorname{Aut}^0(C)&-&\operatorname{Aut}C&-&\unsure{\mathfrak S|\mathcal S}_C&-&1
\end{array}}}
\deleted{$\operatorname{Aut}^0(C)\simeq$}

On a
\diagram{\begin{tikzcd}[ampersand replacement=\&,column sep=small]
1 \arrow[r,dash] \& \operatorname{Aut}^0(C,a) \arrow[r,dash] \& \operatorname{Aut}^0(C) \arrow[r] \& \operatorname{Aut}(\underset{\substack{\Vert\\\operatorname{Aut}_C(a)}}{\pi_a}) \arrow[r] \& 1
\end{tikzcd}}
\deleted{\unsure{À}} Pour déterminer \deleted{$\operatorname{Aut}^0(C,a)=$}
\[
\operatorname{Aut}^0(C,a)=\left\{f\in\operatorname{Aut}(C)\ \middle|\ \operatorname{Ob}f=\operatorname{id}\ .\ (f_a:\operatorname{Aut}(a)-\operatorname{Aut}(\underset{\substack{\Vert\\f(a)}}{a}))=\operatorname{id}\right\}
\]
On considère, pour $\forall x\in C_0$,
\[
T_x=\operatorname{Isom}(a,x),\qquad \pi_a\text{-torseur à droite}
\]
et son commutant \deleted{\unsure{\ensuremath{G_x}}} $\operatorname{Aut}(x)$
